\section{Guida all'Uso e Conduzione di Ulteriori Sperimentazioni}

Questa sezione fornisce le indicazioni operative necessarie alla configurazione dell'ambiente, all'esecuzione degli algoritmi realizzati e all'ispezione delle tabelle prodotte, illustrando inoltre le linee guida per la conduzione di ulteriori sperimentazioni su nuove istanze del problema.

\subsection{Requisiti e Configurazione dell'Ambiente}
Il software sviluppato richiede un interprete Python in versione 3.10 o successiva. Le librerie esterne necessarie all'esecuzione sono elencate nel file \texttt{requirements.txt} e comprendono: \texttt{numpy}, \texttt{pandas} e \texttt{matplotlib}.

Per predisporre l'ambiente di esecuzione isolato, si consiglia la creazione di un ambiente virtuale dedicato tramite i comandi:

\noindent\begin{minipage}{\linewidth}
\begin{verbatim}
# Creazione dell'ambiente virtuale
python3 -m venv .venv

# Attivazione dell'ambiente virtuale
# Su sistemi Unix/Linux e macOS:
source .venv/bin/activate
# Su sistemi Microsoft Windows (Prompt dei comandi):
.\.venv\Scripts\activate.bat
# Su sistemi Microsoft Windows (PowerShell):
.\.venv\Scripts\Activate.ps1

# Installazione delle librerie richieste
pip install -r requirements.txt
\end{verbatim}
\end{minipage}

\subsection{Modalità di Esecuzione}
Il programma principale del progetto è \texttt{demo.py}, che permette sia di eseguire i calcoli su tutti gli scenari generando i risultati, sia di mostrare le tabelle e la rosa ottima. È inoltre disponibile lo script \texttt{main.py} per eseguire i test in modo automatico da terminale.

\subsubsection{Demo e Sperimentazione}
Al fine di garantire la massima portabilità tra sistemi operativi differenti, nella radice del progetto sono stati predisposti specifici script di lancio automatico della demo:
\begin{itemize}
	\item \textbf{Sistemi Unix/Linux e macOS:} tramite lo script Bash \texttt{./run\_demo.sh} (previa assegnazione dei permessi di esecuzione con \texttt{chmod +x run\_demo.sh});
	\item \textbf{Sistemi Microsoft Windows:} tramite lo script batch \texttt{run\_demo.bat} (avviabile da riga di comando o direttamente dall'interfaccia grafica);
	\item \textbf{Invocazione standard Python:} mediante il comando \texttt{python demo.py} (o \texttt{python3 demo.py}).
\end{itemize}

All'avvio, il programma mostra una schermata articolata in due fasi:
\begin{enumerate}
	\item \textbf{Fase di calcolo ed elaborazione degli scenari:} il programma richiede se avviare la sperimentazione completa. Rispondendo affermativamente, il software elabora in sequenza tutti gli undici scenari presenti nella cartella \texttt{scenarios/}, eseguendo per ciascuno le sei varianti algoritmiche su $K = 15$ iterazioni, misurando tempi e picchi di memoria tramite \texttt{tracemalloc}, salvando le metriche in \texttt{results/benchmark\_metrics.csv}, esportando le tabelle dinamiche CSV e generando in automatico i tre grafici PNG di sintesi prestazionale. Qualora i risultati siano già stati calcolati in precedenza, l'utente può saltare tale fase per accedere istantaneamente ai dati già archiviati.
	\item \textbf{Fase di esplorazione interattiva:} conclusa la fase precedente, si apre il menu principale che offre cinque voci:
	\begin{itemize}
		\item \textbf{Opzione 1 (Mostra rosa ottimale):} visualizza in formato tabellare la rosa ottimale identificata dall'algoritmo di ricostruzione per lo scenario selezionato, con i dettagli di costo e rendimento per ciascun ruolo, accompagnata da un riepilogo (WAR totale, spesa sostenuta e budget residuo);
		\item \textbf{Opzione 2 (Visualizza tabelle):} permette di ispezionare a video le matrici $m$ ed $s$. Nelle istanze compatte (come lo scenario demo), le tabelle vengono stampate integralmente, mentre nelle istanze con budget elevato viene mostrata una visualizzazione compatta composta dalle prime 4 colonne, da una colonna di separazione e dalle ultime 4 colonne. In entrambi i casi, viene evidenziato il valore ottimo $m[N, B]$;
		\item \textbf{Opzione 3 (Confronta le varianti algoritmiche):} esegue in tempo reale sullo scenario corrente tutte le sei varianti, stampando una tabella comparativa che certifica la perfetta uguaglianza del valore ottimo calcolato ed esplicita quali varianti supportano la ricostruzione della rosa e quali si limitano al calcolo del solo valore;
		\item \textbf{Opzione 4 (Cambia scenario):} permette di cambiare lo scenario attivo selezionandolo dall'elenco di quelli disponibili;
		\item \textbf{Opzione 5 (Grafici e benchmark delle prestazioni):} rimanda a un sotto-menu che mostra altre voci per la visualizzazione dei grafici generati.
	\end{itemize}
\end{enumerate}

\subsubsection{Ispezione delle Tabelle Generate}
Per la visualizzazione integrale delle tabelle, nei casi in cui il budget assuma dimensioni realistiche rendendo impossibile una lettura a schermo, è possibile consultare i file generati nella cartella \texttt{results/}:
\begin{itemize}
	\item \texttt{matrice\_m\_<scenario>.csv}: contiene la matrice bidimensionale dei valori WAR, con intestazioni di riga associate a ciascun ruolo e intestazioni di colonna relative al budget discretizzato $b \in \{0, \dots, B\}$;
	\item \texttt{matrice\_s\_<scenario>.csv}: contiene la matrice bidimensionale delle scelte, in cui ogni cella $s[i, b]$ riporta l'indice $j \ge 0$ del candidato ingaggiato o il valore convenzionale $-1$ qualora per il ruolo $i$ non sia stato selezionato alcun giocatore.
\end{itemize}

\subsection{Conduzione di Ulteriori Sperimentazioni}
Il software è progettato per consentire la sperimentazione su nuove istanze, variando sia le caratteristiche dimensionali sia i parametri di campionamento.

\subsubsection{Formato dei Dati di Ingresso}
Le nuove istanze del problema devono essere codificate in file con estensione \texttt{.json} e collocate all'interno della cartella \texttt{scenarios/}. La struttura sintattica attesa dal modulo di caricamento (\texttt{src/core/loader.py}) è la seguente:

\noindent\begin{minipage}{\linewidth}
\begin{verbatim}
{
  "scenario_name": "scenario_personalizzato",
  "budget": 500000000,
  "cost_unit": 100000,
  "roles": [
    {
      "role_name": "Pitcher",
      "candidates": [
        { "name": "Player A", "cost": 15000000, "war": 4.5 },
        { "name": "Player B", "cost": 10000000, "war": 3.8 }
      ]
    }
  ]
}
\end{verbatim}
\end{minipage}
Il modulo calcola automaticamente la dimensione dell'istanza e ordina gli scenari in base alla complessità prima dell'esecuzione dei test.

\subsubsection{Linee Guida per Ulteriori Sperimentazioni}
Per estendere la sperimentazione oltre la variazione del budget, il software permette di testare nuove configurazioni modificando semplicemente i parametri nei file JSON:
\begin{itemize}
	\item \textbf{Variazione dei ruoli ($N$):} testando istanze con un numero di ruoli diverso da $N = 9$ (tipico del baseball), è possibile verificare che tempo e memoria crescano in modo lineare, coerentemente con le complessità teoriche $O(N \cdot B \cdot P)$ e $\Theta(N \cdot B)$;
	\item \textbf{Variazione dei candidati ($P$):} modificando il numero di giocatori per ruolo si può confermare che l'occupazione di memoria resta costante (essendo indipendente da $P$) e analizzare l'efficacia del filtro, che scarta una quota percentuale maggiore di giocatori dominati quando il numero di candidati aumenta.
\end{itemize}

\newpage
